%======================================================================
%  3. Data: source, description, processing
%======================================================================
\section{Data}
\label{sec:data}

\subsection{Source}
Data used in the project comes from official public registers maintained by Statistics Denmark (\url{https://www.dst.dk/en}) 
and was retrieved via the public StatBank Denmark database (\url{https://www.statbank.dk}). The dataset combines two main statistical tables with spatial data:
\begin{itemize}
\item \textbf{EJEN77 (Sales of real property by type of transfer, key figures, category of real property, region and time)}: Quarterly dataset recording property sales transactions broken down by region, property category including one-family houses and owner-occupied flats and key indicators, such as number of transactions included in the calculation and average price per property in DKK. Direct link to the table: \url{https://www.statbank.dk/EJEN77}.
\item \textbf{INDKF111 (Income for families by type of income, family type, unit, region and time)}: Yearly dataset reporting family income broken down by region and family type: all families, couples and singles, income category with both disposable income after tax and pre-tax income. Direct link to the table: \url{https://www.statbank.dk/INDKF111}.
\item \textbf{Geospatial data (NUTS-3 / landsdele administrative boundaries)}: The polygons of the 11 Danish provinces (landsdele) used to generate the boundary map were obtained from the official Danish administrative division dataset (DAGI – Danmarks Administrative Geografiske Inddeling, provided by \url{https://data-science.dk/}) and processed into TopoJSON format (landsdele.topo.json). Direct link to the table: \url{https://data-science.dk/datasat/gis/admin/aktuel-admin/danmarks-landsdele-vektor/}.
\end{itemize}
\subsection{Description}

The main dataset (\texttt{affordability\_annual.csv}) covers the years
1995--2024. TODO --- state the exact number of records and any special
feature of the data (missing values, suppressed cells, unbalanced panels,
nominal vs.\ real prices, \ldots).

\begin{table}[H]
\centering
\caption{Variables in the main dataset.}
\label{tab:variables}
\begin{tabularx}{\textwidth}{@{}l l X@{}}
\toprule
\textbf{Variable} & \textbf{Type} & \textbf{Description} \\
\midrule
\texttt{year}          & integer     & Year in which the data was recorded (1995--2024). \\
\texttt{nuts3}         & categorical & NUTS-3 code of the region (e.g.\ \texttt{DK011}). \\
\texttt{landsdel}      & categorical & Danish region / geographical area. \\
\texttt{property\_type}& categorical & Property type: \texttt{one\_family\_house}, \texttt{owner\_occupied\_flat}. \\
\texttt{family\_type}  & categorical & Household type: \texttt{all\_families}, \texttt{couples}, \texttt{singles}. \\
\texttt{n\_sales}      & integer     & Number of property sales recorded. \\
\texttt{avg\_price\_dkk}& numeric    & Average property price in DKK. \\
\texttt{n\_families}   & integer     & Number of families represented. \\
\texttt{avg\_income\_dkk}& numeric   & Average (disposable) income in DKK. \\
\texttt{ratio}         & numeric     & Affordability ratio, \texttt{avg\_price\_dkk} / \texttt{avg\_income\_dkk}. \\
\bottomrule
\end{tabularx}
\end{table}

\noindent In addition, \texttt{landsdele.topo.json} provides the TopoJSON
geometries of the Danish \emph{landsdele}, used for the choropleth map.

\subsection{Data Processing}

The housing and income tables were combined into an annual panel covering the 11 landsdele from 1995 to 2024, the years common to both sources. 
For housing, only ordinary free trades were retained, since family transfers and other transaction types do not reflect market prices. 
One-family houses form the main series and owner-occupied flats a secondary one, as houses alone would misrepresent Byen København, where most purchases are flats.

The housing data are quarterly, whereas income is reported annually. 
Annual prices were therefore obtained as the sales-weighted mean of the quarterly averages,
\begin{equation}
  \bar{p}_{r,t} = \frac{\sum_{q=1}^{4} n_{r,t,q}\, p_{r,t,q}}{\sum_{q=1}^{4} n_{r,t,q}},
\end{equation}
where $p_{r,t,q}$ is the average price and $n_{r,t,q}$ the number of sales in the price calculation for landsdel $r$ in quarter $q$ of year $t$. 
This reproduces the average over all sales in the year. 
Quarters without sales, which the source reports as a price of zero, were treated as missing and given no weight. 
As prices are averages rather than medians, the mix of homes sold may change over time.

Income was taken as the average income per family, both after tax (disposable income) and before tax (pre-tax income), for all families, couples and single people. 
The published totals were used as reported. 
Affordability is expressed as the price-to-income ratio
\begin{equation}
  R^{\,k}_{r,t} = \frac{\bar{p}_{r,t}}{\bar{y}^{\,k}_{r,t}}, \qquad k \in \{\text{disposable},\ \text{pre-tax}\},
\end{equation}
where $\bar{y}^{\,k}_{r,t}$ is the average family income of type $k$. 
Both quantities are in current prices, so the ratio requires no inflation adjustment. 
For couples and single people, the same average price is divided by the respective group's income, regardless of who actually bought the property.

The tables were joined on the NUTS~3 code. 
Region labels in the Statistics Denmark exports were harmonised (the ``Province'' prefix removed and Unicode forms normalised) and matched to the names in the boundary file, which carries the corresponding codes. 
All 11 landsdele matched, and the resulting panel is complete, with no missing values.

The boundary file was prepared for web display in three steps. 
Islets smaller than $0.5~\text{km}^2$, including Ertholmene, were removed. 
The geometry was then simplified using weighted Visvalingam simplification, keeping 5\% of the points while preserving all shapes, which reduced the file from 2.8~MB to about 180~KB. 
The result was stored as TopoJSON with true geographic coordinates. 
Bornholm is not repositioned in the file and will be drawn as an inset at the visualisation stage.
